

\subsection*{Resume the annual book from its last committed operation}

A saved balance sheet is not enough to restart this simulation.  After the
first inventory sale, the remaining unit still carries its acquisition identity,
while each tax payable belongs to an assessment and retains any settlements.
The next operation must also be known: repeating the sale would change both
stock and cash even if the repeated journal balanced.  The optional
\texttt{Annual\allowbreak{}Inventory\allowbreak{}Store/1} therefore preserves the
complete annual request, the full period book and the ordered operation history.

The local store records the evaluated opening before attempting annual
operations.  Each later transaction records one attempted operation together
with its complete result and its position in the original request.  A successful
sale advances both the attempt position and the successful transaction history.
A refused sale advances the attempt position but leaves the last valid financial
state unchanged.  The refusal and all remaining operations stay visible, and
that immutable request stops.  Repairing its assumptions creates a different
request rather than rewriting the recorded history.

An explicit operation budget permits work to pause between transactions.
While operations remain unattempted, the stored result says that it is paused;
it does not turn a valid partial-year statement into a completed run.  A caller
can also supply the position it expected before a batch.  If that batch has
already committed, a retry returns the checked current result without applying
another sale.  If the caller expects a position that has not been reached, the
store rejects the request.

In the running example, the first sale leaves cash of 467, one recovered unit
costing 10, and the carried tax payable of \(41/4\).  Closing and reopening the
store retains all 26 accounts and their associated histories.  Continuing through
the declared tax assessment and period carry preserves the residual unit; its
sale in the following year gives cash of 482 and a payable of \(49/4\), exactly
as uninterrupted execution does.  The complete statements and refused prefixes,
not merely these selected balances, are compared in the reference checks.

Reopening checks the saved request, implementation identities, row order and
content hashes, then recomputes each committed prefix through the same financial
evaluator.  This catches altered balances even when a stored hash has also
been changed.  Replayed calculation has no external financial side effects.
A new operation uses the loaded full state directly; reopening deliberately
pays the cost of replay to validate that state.

The implementation uses one local SQLite writer and an atomic transaction for
each operation result.  In the process-exit checks, termination before commit
leaves the preceding result, while termination after commit retains exactly one
new result.  A fresh process can then reopen and continue.  These checks support
local restart behavior for this finite annual profile.  They do not establish
hardware power-loss durability, concurrent-writer correctness, efficient scaling
or authenticity against an attacker who controls the entire database.  Family
scheduling, migration and broader execution profiles remain separate work.
